\renewcommand{\chaptertitle}{A Virtual Machine Runtime}

%  Make first page footer:
\fancypagestyle{firststyle}{%
\fancyhf{}% Clear header/footer
\fancyhead{}
\fancyfoot[L]{{\footnotesize
              %% We toggle between these:
              % Manuscript submitted for review.\\
              \authorname{}~\authorpatp{}.  ``\chaptertitle{}''.  {\it Nockonomicon} (2025):  50–$x$.
              }}
}
%  Arrange subsequent pages:
\fancyhf{}
\fancyhead[LE]{{\urbitfont Nockonomicon}}
\fancyhead[RO]{\chaptertitle{}}
\fancyfoot[LE,RO]{\thepage}

\chapter{\chaptertitle}

\thispagestyle{firststyle}

\begin{quote}
\textbf{The Kraken}
Below the thunders of the upper deep,
Far, far beneath in the abysmal sea,
His ancient, dreamless, uninvaded sleep
The Kraken sleepeth: faintest sunlights flee
About his shadowy sides; above him swell
Huge sponges of millennial growth and height;
And far away into the sickly light,
From many a wondrous grot and secret cell
Unnumbered and enormous polypi
Winnow with giant arms the slumbering green.
There hath he lain for ages, and will lie
Battening upon huge sea worms in his sleep,
Until the latter fire shall heat the deep;
Then once by man and angels to be seen,
In roaring he shall rise and on the surface die.
(Alfred, Lord Tennyson)
\end{quote}

To date, we have run Nock either by hand or using a simple interpreter.  While this suffices for relatively small programs, our objective of building a Nock computer requires an ``operating function'', or persistent event loop.  This chapter introduces the considerations of running Nock continuously and reactively as well as how to maintain state and respond to opcode 11 hints.

\section{The Event Loop}

A single Nock expression as a subject--formula pair accepts two nouns and yields a new noun.  We can think of this new noun as the total state of a program or an operating system, feeding it into the next expression as the subject for a new formula to evaluate against.  This preserves the determinism and legibility of Nock while allowing for complex state transitions.

\subsection{A Simple State Machine}

A state machine is a mathematical abstraction for a system that has a definite state at any particular time, and which changes in response to inputs.

\subsubsection{Example:  Stoplight}

Think of a stoplight, which has a state of red, yellow, or green, and which changes its state in response to an external timer or a button press.


\subsubsection{Example:  Revolver}

Or again, of a six-shot revolver, which if we include chamber order has 64 finite state permutations with thirteen possible inputs---to ``load'' or ``clear'' at each chamber and to ``fire''---and two possible outputs---a ``bang'' or a ``click'' accompanied by a new state.

% 6 1
% 5 6
% 4 15
% 3 20
% 2 15
% 1 6
% 0 1

Aggregate or macro actions are also available---a ``moon clip'' corresponds to loading all six chambers at once, and a ``dump cylinder'' corresponds to clearing all six chambers at once.  The revolver is a state machine with a finite number of states, inputs, and outputs.

\footnote{\textct{+poke} and \textct{+peek} derive from \href{https://en.wikipedia.org/wiki/PEEK_and_POKE}{a very old convention in \textct{BASIC}}, wherein \textct{POKE} and \textct{PEEK} commands allowed programmers to write to and read from arbitrary memory addresses.}

We will call a state machine which adheres to the definition of its lifecycle function a \emph{kernel}.  Thus a kernel is a state machine of a certain shape which accepts nouns of a certain shape as input and produces nouns of a certain shape as output.  The kernel is the event loop and some wrapper of library functionality to bound and define the complete behavior of the state machine.  A kernel may be batch or continuous depending on whether its lifecycle function is reentrant.

\subsection{An Evolved State Machine}

While the two-armed model provides a straightforward way to manage state transitions, Nockhounds have found in practice that it is more convenient to supply four arms at the following axes:

\begin{enumerate}
  \item[\textct{+4}:] \textct{+load} is used in kernel upgrades, allowing the event loop to update itself live.
  \item[\textct{+10}:] \textct{+wish} accepts a core and parses it against a standard library for the kernel.
  \item[\textct{+22}:] \textct{+peek} grants read-only access to the system's state; this is often called a ``scry''.
  \item[\textct{+23}:] \textct{+poke} accepts events and processes them, resulting in a new state noun.
\end{enumerate}
